\chapter{Planificación del proyecto}
\label{chap:planificacion}

\ph{Este capítulo demuestra que el TFG no es solo un producto software, sino un
proyecto gestionado de principio a fin: se planifica al inicio (acta de constitución,
interesados, metodología como decisión de planificación, EDT, cronograma como línea
base y presupuesto como línea base de coste), se controla durante la ejecución
(comunicaciones, riesgos, control de cambios, calidad) y se cierra formalmente con
el análisis de desviaciones ---de calendario, esfuerzo y coste--- y las lecciones
aprendidas. Trata el plan como un documento vivo: la estimación inicial no se borra
cuando la realidad la contradice, se compara con ella.}

\section{Acta de constitución del proyecto}
\label{sec:acta_constitucion}

\ph{Resume en media página el «contrato» inicial del proyecto: qué se va a hacer,
bajo qué supuestos y con qué restricciones. No repitas la introducción: referencia
los objetivos (Sección~\ref{sec:objetivos}) y el alcance (Sección~\ref{sec:alcance})
y céntrate en lo que aún no está escrito (supuestos, restricciones, criterio de cierre).}

\begin{table}[H]
\centering
\begin{tabularx}{\textwidth}{p{4cm}X}
\toprule
\textbf{Elemento} & \textbf{Descripción} \\
\midrule

Propósito &
Desarrollar una herramienta que permita generar automáticamente modelos de características en formato UVL para facilitar la creación de modelos de variabilidad software, proporcionando a investigadores y desarrolladores un mecanismo configurable para definir las características, restricciones y atributos de los modelos generados. \\

\midrule

Objetivos &
Los objetivos del proyecto son los definidos en la Sección~\ref{sec:objetivos}. De forma general, el proyecto busca diseñar e implementar un generador automático de modelos de características, integrarlo en UVLHub y proporcionar soporte para diferentes niveles de expresividad del lenguaje UVL. \\

\midrule

Alcance &
El proyecto comprende el diseño e implementación del generador automático de modelos de características, su integración dentro de la plataforma UVLHub y la incorporación progresiva de soporte para modelos booleanos, aritméticos y con tipos definidos por UVL. Quedan fuera del alcance la modificación del funcionamiento interno de Flamapy y el desarrollo de nuevos algoritmos de análisis o razonamiento. \\

\midrule

Supuestos &
Se asume la disponibilidad del ecosistema software necesario para el desarrollo, incluyendo acceso al código fuente de UVLHub y Flamapy, así como las herramientas necesarias para la generación y análisis de modelos de variabilidad. También se considera disponible el entorno de desarrollo necesario para implementar, probar y validar la solución. \\

\midrule

Restricciones &
El proyecto está condicionado por las tecnologías y arquitectura existentes en UVLHub y Flamapy, especialmente el lenguaje Python (su framework Flask) y el estándar UVL. Además, debe desarrollarse antes del 7 de octubre, la fecha límite para la entrega de la memoria. \\

\midrule

Entregables principales &
Los principales entregables del proyecto son el código fuente del generador automático integrado en UVLHub y la memoria del TFG. \\

\midrule

Criterio de cierre &
El proyecto se considerará finalizado cuando el generador permita crear modelos de características funcionales en los niveles de expresividad definidos, se encuentre integrado dentro de UVLHub y se hayan realizado pruebas que validen el correcto funcionamiento tanto de la generación de modelos como del front-end desarrollado para utilizarlo en UVLHub \\

\bottomrule
\end{tabularx}
\caption{Acta de constitución del proyecto.}
\label{tab:acta_constitucion}
\end{table}

\subsection{Interesados del proyecto}
\label{sec:interesados}

\ph{Identifica a las personas con interés en el proyecto y qué espera cada una.
En un TFG el listado es corto; lo importante es que ninguna fila quede sin
expectativa concreta.}

\begin{longtable}{@{}P{3cm} P{3.2cm} P{5.5cm} P{2.5cm}@{}}
  \toprule
  \textbf{Interesado} & \textbf{Rol} & \textbf{Interés / expectativa} & \textbf{Implicación} \\
  \midrule
  \endhead

  Autor del proyecto &
  Responsable del desarrollo &
  Completar el Trabajo Fin de Grado mediante el desarrollo de una solución funcional, documentada y alineada con los objetivos establecidos. &
  Diaria \\[0.3em]

  Tutor del TFG &
  Supervisión y validación académica &
  Garantizar el correcto desarrollo del proyecto, proporcionando orientación técnica y revisando los avances y entregables realizados. &
  Periódica \\[0.3em]

  Grupo de investigación / ecosistema UVLHub-Flamapy &
  Entorno receptor de la solución &
  Disponer de una herramienta que permita generar modelos de características de forma automática y ampliar las capacidades existentes del ecosistema de variabilidad software. &
  Puntual \\[0.3em]

  Tribunal del TFG &
  Evaluación final del proyecto &
  Valorar la calidad de la memoria, la defensa realizada y los resultados obtenidos durante el desarrollo del trabajo. &
  Al cierre \\

  \bottomrule
  \caption{Registro de interesados del proyecto.}
  \label{tab:interesados}
\end{longtable}

\section{Metodología}
\label{sec:metodologia}

\subsection{Metodología de desarrollo}
\label{sec:metodologia_desarrollo}

\ph{Justifica la metodología (o metodologías) elegida(s) para el desarrollo y para
la redacción de la memoria. Es habitual usar enfoques distintos para cada fase.}

\ph{Puntos a cubrir:
\begin{itemize}
  \item \ph{Metodología elegida (Waterfall, Scrum, Kanban, iterativa\dots) y por qué.}
  \item \ph{Qué alternativas consideraste y por qué las descartaste (una frase por alternativa basta).}
  \item \ph{Cómo se gestiona el control de versiones (Git, GitHub/GitLab\dots).}
  \item \ph{Si se usan herramientas de seguimiento (issues, tableros, CI), nómbralas brevemente.}
\end{itemize}}

Para el desarrollo del proyecto se ha seguido una metodología incremental e iterativa, adaptada a la naturaleza evolutiva del generador automático de modelos de características. Debido al carácter investigador y evolutivo del proyecto, la especificación inicial se refinó progresivamente durante el desarrollo conforme se analizaban las necesidades del generador y su integración dentro del ecosistema existente. \cite{sundermann2024generating}.

En una primera etapa se desarrolló una versión inicial del generador orientada a la creación de modelos de características booleanos, junto con su integración dentro de UVLHub mediante la interfaz gráfica proporcionada para configurar el proceso de generación. Una vez obtenida esta primera versión funcional, se realizaron ciclos sucesivos de implementación, validación y corrección que permitieron mejorar la estabilidad del sistema e incorporar nuevas funcionalidades pertinentes a los siguientes niveles de expresividad que se pretendió conseguir. Esta forma de hacerlo permitió validar cada nueva incorporación antes de continuar con las siguientes etapas y reducir el riesgo de introducir errores en funcionalidades previamente implementadas que se pudieran arrastrar a las siguientes.


\subsubsection*{Alternativas consideradas}

No se optó por metodologías secuenciales como el modelo en cascada, ya que la naturaleza del proyecto requería realizar ajustes continuos conforme avanzaba el desarrollo y se identificaban nuevas necesidades relacionadas con la generación de modelos de características. Tampoco se aplicó una metodología Scrum completa, puesto que el desarrollo no se organizó mediante sprints, roles o hitos propios de este marco de trabajo.


\subsubsection*{Control de versiones}

El control de versiones del proyecto se ha realizado mediante Git, utilizando
GitHub como plataforma de almacenamiento y seguimiento de la evolución del
código. Debido a la evolución del proyecto y a la necesidad de integrar el
generador dentro de diferentes versiones de UVLHub, el desarrollo quedó
distribuido entre varios repositorios y una solicitud de incorporación de
cambios (\textit{pull request}) en el repositorio oficial del proyecto.

Por este motivo, el repositorio
\href{https://github.com/antjimort/uvlhub}{\texttt{antjimort/uvlhub}}
no representa por sí solo todo el historial del Trabajo Fin de Grado. Dicho repositorio recoge principalmente la fase final de integración,
refactorización, corrección y validación realizada entre mayo y septiembre de
2026. La actividad desarrollada entre mayo y julio corresponde a la fase
principal de estabilización, mientras que entre el 14 y el 18 de septiembre se
realizaron correcciones adicionales derivadas de la revisión técnica del tutor.

\begin{table}[H]
\centering
\small
\renewcommand{\arraystretch}{1.15}

\begin{tabularx}{\textwidth}{@{}
    >{\raggedright\arraybackslash}p{0.25\textwidth}
    >{\centering\arraybackslash}p{0.19\textwidth}
    >{\raggedright\arraybackslash}X
@{}}
\toprule

\textbf{Recurso} &
\textbf{Periodo relevante} &
\textbf{Papel dentro del desarrollo} \\

\midrule

\href{https://github.com/antjimort/fm_generator}
{\texttt{fm\_generator}} &
15/05/2025--21/04/2026 &
Repositorio inicial del componente generador. Contiene 38 commits realizados
por el autor relacionados con la implementación de niveles UVL, atributos,
restricciones, cardinalidades y opciones de configuración. Los commits
anteriores a mayo de 2025 no se consideran parte de la contribución del TFG. \\[0.4em]

\href{https://github.com/antjimort/uvlhub_v2}
{\texttt{uvlhub\_v2}} &
14/02/2026--03/04/2026 &
Repositorio transitorio utilizado durante la migración de los cambios previos
a una nueva versión de UVLHub. En él se realizaron ajustes de migración,
persistencia de parámetros del asistente y avances en los niveles booleano y
aritmético del generador. \\[0.4em]

\href{https://github.com/diverso-lab/uvlhub/pull/117}
{\texttt{PR \#117}} &
06/04/2026--19/04/2026 &
Solicitud de incorporación de cambios en el repositorio oficial de UVLHub,
titulada \textit{feat: boolean level ready}. La PR contiene 46 commits en
total, de los cuales 45 fueron realizados por el autor, y fue fusionada en la
rama \texttt{generator} el 19 de abril de 2026. \\[0.4em]

\href{https://github.com/antjimort/uvlhub}
{\texttt{uvlhub}} &
05/2026--09/2026 &
Repositorio utilizado durante la fase final de integración, estabilización y
consolidación técnica. Incluye refactorizaciones de la arquitectura,
correcciones de compatibilidad, adaptación de pruebas, automatización de
despliegue y correcciones finales relacionadas con la validación de modelos
UVL, PySAT y la integración del generador como \textit{plugin} de Flamapy. \\

\bottomrule
\end{tabularx}

\caption{Trazabilidad de los repositorios y recursos utilizados durante el desarrollo.}
\label{tab:trazabilidad-repositorios}
\end{table}

Los repositorios y la solicitud de incorporación de cambios anteriores no
deben interpretarse como desarrollos independientes que puedan sumarse entre
sí, ya que parte de los commits realizados en el repositorio
\texttt{fm\_generator} fueron posteriormente incorporados en el PR~\#117.
La Tabla~\ref{tab:trazabilidad-repositorios} permite identificar el origen de
cada fase del desarrollo sin duplicar la dedicación contabilizada en el
cronograma.

La convención \textit{Conventional Commits} se adoptó progresivamente durante
el proyecto. Por este motivo, algunos commits iniciales no siguen una
nomenclatura completamente normalizada. Sin embargo, en las fases posteriores
se utilizaron prefijos como \textit{feat} para la incorporación de nuevas
funcionalidades, \textit{fix} para la corrección de errores, \textit{refactor}
para reorganizaciones internas del código, \textit{test} para cambios en las
pruebas, \textit{chore} para tareas de mantenimiento, \textit{style} para
ajustes de estilo y \textit{ci} para modificaciones de integración continua.

Esta organización permite mantener la trazabilidad de la evolución completa
del generador, desde la implementación inicial como componente independiente
hasta su integración y estabilización dentro de UVLHub.

\subsubsection*{Herramientas de seguimiento}

El seguimiento del desarrollo se ha realizado principalmente mediante el propio historial de GitHub, utilizando los commits realizados como registro de la evolución del proyecto y de las decisiones técnicas tomadas durante la implementación.

Además, las reuniones periódicas con el tutor han permitido revisar el estado del desarrollo para analizar las deficiencias encontradas con el propósito de corregirlas y llevar la calidad del generador al siguiente nivel. Los detalles de estas reuniones se detallarán en el siguiente apartado.


\subsection{Metodología de validación y evaluación}
\label{sec:metodologia_validacion}

\ph{Explica cómo se medirá que el proyecto cumple lo prometido. Normalmente combina
dos vías: (a) pruebas y trazabilidad requisitos--pruebas (Capítulo~\ref{chap:pruebas}),
y (b) evaluación de resultados frente a los criterios de éxito de la
Tabla~\ref{tab:criterios-exito}, cuyo análisis se presenta en el
Capítulo~\ref{chap:resultados}. Aquí basta con declarar el método; los datos van
en sus capítulos.}

La validación del proyecto se ha realizado mediante un proceso iterativo durante el desarrollo, combinando pruebas funcionales, pruebas automatizadas y revisiones periódicas con el tutor del proyecto.

Tras la implementación de cada nueva funcionalidad, se realizaron comprobaciones manuales para verificar que el generador producía modelos de características correctos según los parámetros configurados, incluyendo los diferentes niveles de expresividad soportados por UVL, así como la correcta interacción del usuario con el front-end desarrollado en UVLHub. Además, se ejecutaron pruebas unitarias automatizadas para detectar posibles errores o regresiones, incorporando nuevos casos de prueba únicamente cuando aportaban valor al sistema; es decir, que encontraran posibles bugs reales y relevantes.

La evaluación del cumplimiento del proyecto se ha realizado siguiendo los criterios de éxito definidos en la Tabla~\ref{tab:criterios-exito}. Estos criterios permiten comprobar de forma objetiva aspectos como la generación automática de modelos UVL, la configuración del proceso de generación, el soporte de construcciones avanzadas del lenguaje, la integración dentro del ecosistema Flamapy y UVLHub, y la existencia de una batería de pruebas automatizadas.

Además de la validación técnica continua, se realizaron revisiones del estado del proyecto junto con el tutor. Cuando se alcanzó aproximadamente el 80\% de desarrollo, se realizó una primera revisión en la que se concluyó que la solución era funcionalmente válida, pero presentaba ciertas deficiencias técnicas que debían ser corregidas antes de considerarla finalizada. Entre ellas se identificaron problemas relacionados con la integración del generador dentro de UVLHub, la pérdida de algunas funcionalidades existentes durante determinadas modificaciones (que posteriormente se recuperaron sin problema), la reducción de la cobertura de pruebas, la presencia de elementos no deseados en el repositorio y aspectos de calidad del código (documentación y limpieza de la implementación).

Tras aplicar las correcciones necesarias, se realizó una segunda revisión con el tutor para comprobar el estado de la nueva versión. En esta ocasión únicamente se detectaron ajustes menores, que fueron solucionados posteriormente, obteniendo finalmente la validación positiva del tutor.

\section{Estructura de descomposición del trabajo (EDT)}
\label{sec:edt}

\ph{Desglosa el proyecto en paquetes de trabajo con identificador jerárquico de dos
niveles. La regla de oro: todo entregable debe poder marcarse como hecho/no hecho
sin ambigüedad. Los IDs se reutilizan en el cronograma y en el cierre para dar
trazabilidad. Con 6--10 paquetes es suficiente. Desglosa al segundo nivel solo
donde aporte; los paquetes hoja deben corresponder a entregables.}

Este proyecto se podría desglosar de la siguiente manera:

\begin{longtable}{@{}P{1cm} P{4.5cm} P{8.8cm}@{}}
  \toprule
  \textbf{ID} & \textbf{Paquete de trabajo} & \textbf{Entregable verificable} \\
  \midrule
  \endhead

  1 & Gestión y planificación del proyecto &
  Planificación inicial del proyecto, estudio de la solución propuesta por los investigadores
  a partir del trabajo de referencia y seguimiento del desarrollo realizado. \\[0.3em]

  1.1 & Planificación inicial y estudio de la solución &
  Acta de constitución, definición del alcance, planificación temporal del proyecto y
  análisis de la solución propuesta en el artículo de referencia. \\[0.3em]

  1.2 & Seguimiento y control &
  Reuniones de seguimiento, revisión del estado del desarrollo y gestión de los cambios
  realizados durante el proyecto. \\[0.3em]

  2 & Análisis y diseño de la solución &
  Diseño de la arquitectura del generador y definición de los componentes necesarios para su implementación. \\[0.3em]

  2.1 & Análisis de requisitos &
  Catálogo de requisitos funcionales y técnicos del generador automático de modelos. \\[0.3em]

  2.2 & Diseño del generador y su integración &
  Arquitectura definida para la generación de modelos, configuración del sistema e integración con UVLHub. \\[0.3em]

  3 & Desarrollo del generador automático &
  Implementación del generador capaz de producir modelos de características en formato UVL. \\[0.3em]

  3.1 & Implementación de niveles de expresividad &
  Soporte implementado para generación de modelos booleanos, aritméticos y con tipos definidos por UVL. \\[0.3em]

  3.2 & Desarrollo de configuración e interfaz &
  Sistema de configuración del generador e integración del flujo de interacción mediante el front-end de UVLHub. \\[0.3em]

  4 & Pruebas y validación &
  Suite de pruebas automatizadas y validación funcional del generador integrado. \\[0.3em]

  5 & Documentación y defensa &
  Memoria final del Trabajo Fin de Grado y material necesario para la presentación de defensa. \\

  \bottomrule

  \caption{Estructura de descomposición del trabajo (EDT).}
  \label{tab:edt}
\end{longtable}

\section{Cronograma e hitos (línea base)}
\label{sec:fases}

\ph{Divide el proyecto en fases claras con un hito (entregable verificable) por fase.
La línea base tiene DOS dimensiones: el esfuerzo (columnas \textbf{Est.}\ y
\textbf{Real}, en horas) y el calendario (columnas \textbf{Fin prev.}\ y
\textbf{Fin real}, fechas de hito en formato mm/aa). Ambas se fijan al inicio y NO
se reescriben cuando la realidad cambia: la estimación inicial se conserva y se
compara con lo real, y los cambios significativos se registran en la
Sección~\ref{sec:control_cambios}. La columna \textbf{EDT} indica qué paquetes de la
Tabla~\ref{tab:edt} cubre cada fase, y la columna \textbf{Obj.}\ qué objetivo técnico o formativo relacionado aborda. En particular, las
fases se vinculan principalmente con los objetivos O1--O6 definidos en la
Sección~\ref{sec:objetivos}.}

\ph{Tabla de fases del desarrollo: rellena las filas con tus propias fases.}

Teniendo como fecha de inicio del proyecto el 12 de abril de 2025, el
cronograma real quedó dividido en varias etapas de desarrollo, integración y
validación. La Tabla~\ref{tab:cronograma-ejecucion} recoge los periodos de
actividad identificables a partir del historial de los repositorios y de los
hitos de integración del proyecto.

\begin{table}[H]
\centering
\small
\setlength{\tabcolsep}{3pt}
\renewcommand{\arraystretch}{1.2}

\begin{tabularx}{\textwidth}{@{}
    >{\centering\arraybackslash}p{1.8cm}
    >{\centering\arraybackslash}p{1.5cm}
    >{\centering\arraybackslash}p{1.2cm}
    >{\raggedright\arraybackslash}X
    >{\raggedright\arraybackslash}X
@{}}
\toprule

\shortstack[c]{\textbf{Periodo}\\\textbf{real}} &
\shortstack[c]{\textbf{Fin previsto}\\\textbf{del paquete}} &
\textbf{EDT} &
\textbf{Actividad principal} &
\textbf{Hito o resultado} \\

\midrule

04/2025 &
04/2025 &
1.1 &
Análisis inicial y planificación del proyecto. &
Planificación inicial y estudio de viabilidad del generador. \\[0.3em]

04/2025 &
04/2025 &
2.1--2.2 &
Análisis de requisitos y diseño de la solución. &
Requisitos definidos y prototipo inicial preparado para validar la integración. \\[0.3em]

05/2025--07/2025 &
05/2026 &
3 &
Primera iteración del generador y de su integración inicial con UVLHub. &
Implementación de restricciones booleanas, cardinalidades, atributos manuales,
configuración inicial del \textit{wizard} e integración preliminar del componente. \\[0.3em]

08/2025--01/2026 &
--- &
--- &
Periodo de inactividad del proyecto. &
No se realizaron actividades ni se establecieron hitos durante este intervalo. \\[0.3em]

02/2026--04/2026 &
05/2026 &
3 &
Reanudación del desarrollo, migración a UVLHub 2.0 e incorporación de nuevos
niveles de expresividad. &
Migración funcional, soporte del nivel booleano, funciones agregadas e
integración mediante el PR~\#117, fusionado el 19 de abril de 2026. \\[0.3em]

05/2026--07/2026 &
05/2026 &
3--4 &
Integración final, refactorización, adaptación de pruebas y validación del
sistema. &
Generador estabilizado, pruebas adaptadas a la nueva arquitectura y validación
final del tutor. \\

14/09/2026--18/09/2026 &
--- &
3--4 &
Correcciones finales y consolidación técnica posterior a la fase principal de
validación. &
Corrección de la generación de tipos y expresiones, validación de modelos UVL
con \texttt{UVLReader} y PySAT, integración del generador como \textit{plugin}
de Flamapy, reorganización de las pruebas unitarias y limpieza de elementos no
deseados del repositorio. \\

\bottomrule
\end{tabularx}

\caption{Cronograma de ejecución e hitos del proyecto.}
\label{tab:cronograma-ejecucion}
\end{table}

La distribución del esfuerzo se mantiene agrupada por paquetes de trabajo,
ya que las horas reales se registraron por fases generales y no de forma
independiente para cada una de las iteraciones temporales anteriores. Esta
decisión evita asignar artificialmente horas concretas a periodos para los que
no existe un registro horario detallado.

\begin{table}[H]
\centering
\small
\setlength{\tabcolsep}{3pt}
\renewcommand{\arraystretch}{1.2}

\begin{tabularx}{\textwidth}{@{}
    >{\centering\arraybackslash}p{0.55cm}
    >{\raggedright\arraybackslash}X
    >{\centering\arraybackslash}p{1.35cm}
    >{\centering\arraybackslash}p{1.25cm}
    >{\centering\arraybackslash}p{1.25cm}
    >{\centering\arraybackslash}p{1.25cm}
    >{\raggedright\arraybackslash}X
@{}}
\toprule

\textbf{\#} &
\textbf{Fase} &
\textbf{EDT} &
\textbf{Obj.} &
\shortstack[c]{\textbf{Est.}\\\textbf{(h)}} &
\shortstack[c]{\textbf{Real}\\\textbf{(h)}} &
\textbf{Hito / entregable} \\

\midrule

1 &
Análisis inicial y planificación del proyecto &
1.1 &
O1 &
12 &
12 &
Prototipo inicial y planificación del desarrollo validada. \\[0.3em]

2 &
Análisis y diseño de la solución &
2.1--2.2 &
O1--O2 &
25 &
25 &
Requisitos definidos, análisis de la solución propuesta y prototipo inicial
desarrollado para validar la viabilidad de la integración. \\[0.3em]

3 &
Desarrollo e integración del generador automático &
3 &
O1--O4 &
593 &
703 &
Generador implementado e integrado dentro de UVLHub con soporte para los
niveles de expresividad definidos y corrección de errores detectados durante
el desarrollo. \\[0.3em]

4 &
Pruebas y validación de la solución &
4 &
O5 y O6 &
120 &
170 &
Suite de pruebas automatizadas, validaciones funcionales y revisiones finales
del sistema completadas. \\

\midrule

&
\textbf{Total} &
&
&
\textbf{750} &
\textbf{910} &
\\

\bottomrule
\end{tabularx}

\caption{Fases de desarrollo y esfuerzo asociado: línea base frente a valores reales.}
\label{tab:fases-desarrollo}
\end{table}

La fase principal de desarrollo, integración y validación concluyó en julio
de 2026. Posteriormente, entre el 14 y el 18 de septiembre de 2026, se
realizó una fase adicional de correcciones y consolidación técnica como
respuesta a las observaciones del tutor. Por tanto, el cierre técnico
definitivo del proyecto se sitúa en septiembre de 2026.

La línea base preveía finalizar las fases principales en mayo de 2026. La
finalización de la fase principal en julio supuso una desviación de dos meses,
mientras que el cierre técnico definitivo en septiembre representó una
desviación acumulada de cuatro meses respecto a la previsión inicial.

En términos de esfuerzo, las 910 horas reales de desarrollo superaron en
160 horas las 750 horas inicialmente estimadas, lo que representa una
desviación del 21,3\,\%. El esfuerzo adicional se concentró en las fases de
desarrollo e integración del generador y de pruebas y validación, debido a
los problemas de compatibilidad con dependencias externas, la necesidad de
realizar refactorizaciones, la adaptación de las pruebas y las correcciones
finales realizadas en septiembre tras la revisión del tutor.

De las 910 horas de desarrollo, 400 corresponden al contrato de apoyo
técnico asociado a TASOVA Plus y las 510 restantes son atribuibles al
desarrollo del TFG. Junto con las 70 horas dedicadas a la redacción,
maquetación y revisión de la memoria, la dedicación atribuible al TFG
asciende a 580 horas.

% --- Opcional: diagrama de Gantt (espacio reservado; no requiere cargar paquetes) ---
% \begin{figure}[htbp]
%   \centering
%   \includegraphics[width=\textwidth]{figures/gantt}
%   \caption{Diagrama de Gantt del proyecto (línea base de calendario).}
%   \label{fig:gantt}
% \end{figure}

\ph{Opcional: incluye un diagrama de Gantt como figura para visualizar la línea
base de calendario y sus desviaciones; exporta el diagrama como imagen y
descomenta el bloque anterior.}

\ph{Tabla equivalente para la redacción de la memoria, agrupada en bloques de
capítulos con hitos distintos y verificables; el desglose por capítulo individual
es opcional. Todas las filas corresponden al paquete 5 de la EDT.}


\begin{longtable}{@{}P{1cm} P{4cm} c c c P{5cm}@{}}
\toprule
\textbf{\#} & \textbf{Capítulos} & \textbf{EDT} & \textbf{Est.} & \textbf{Real} & \textbf{Hito / Entregable} \\
\midrule
\endhead

1 &
Caps. 1--3: introducción, estado del arte y planificación
& 5
& 30 h
& 33 h
& Borrador de contexto y plan aprobado por el tutor. \\[0.3em]

2 &
Caps. 4--7: núcleo técnico (requisitos, diseño, implementación y pruebas)
& 5
& 25 h
& 24 h
& Núcleo técnico redactado y trazable con el repositorio. \\[0.3em]

3 &
Caps. 8--9: resultados y conclusiones
& 5
& 15 h
& 13 h
& Memoria completa lista para revisión final y depósito. \\

\midrule
\textbf{Total}
&
&
&
\textbf{70 h}
&
\textbf{70 h}
&
\\

\bottomrule
\caption{Fases de redacción de la memoria (estimadas vs. reales).}
\label{tab:fases-redaccion-memoria}
\end{longtable}

Las 70 horas reales indicadas en la
Tabla~\ref{tab:fases-redaccion-memoria} corresponden al trabajo completo de
elaboración de la memoria, incluyendo la redacción y revisión de los capítulos,
la organización de referencias, la preparación de tablas y figuras, la
maquetación en \LaTeX, la revisión de referencias cruzadas y la preparación de
los apéndices.

La distribución de 33, 24 y 13 horas corresponde a bloques de capítulos y no
a una estimación basada únicamente en su extensión. Las tareas técnicas
necesarias para obtener resultados, ejecutar pruebas y recopilar evidencias se
contabilizaron dentro de las fases de desarrollo y validación, evitando sumar
dos veces el mismo esfuerzo.

\ph{Tabla de resumen global: agrega las dos tablas anteriores. Los totales DEBEN
coincidir con los de las tablas de origen (Tablas~\ref{tab:fases-desarrollo}
y~\ref{tab:fases-redaccion-memoria}).}

\begin{table}[H]
\centering
\begin{tabular}{@{}P{4cm} c c@{}}
\toprule
\textbf{\# Fase} & \textbf{Est.} & \textbf{Real} \\
\midrule
Desarrollo (Tabla~\ref{tab:fases-desarrollo})
& 750 h & 910 h \\[0.3em]

Memoria (Tabla~\ref{tab:fases-redaccion-memoria})
& 70 h & 70 h \\

\midrule
\textbf{Total}
& \textbf{820 h} & \textbf{980 h} \\

\bottomrule
\end{tabular}
\caption{Resumen global del proyecto (estimadas vs. reales).}
\label{tab:resumen-global-proyecto}
\end{table}

El total de 980 horas representa la dedicación real acumulada durante el
proyecto, pero no constituye íntegramente la carga académica del TFG. De este
total, 400 horas corresponden al contrato de apoyo técnico asociado a TASOVA
Plus y no forman parte de la dedicación académica del TFG. Las 580 horas
restantes sí son atribuibles al TFG, incluyendo 510 horas de desarrollo y
70 horas dedicadas a la redacción, maquetación y revisión de la memoria.

Por tanto, la comparación con las 300 horas asociadas a un TFG de 12 ECTS
debe realizarse utilizando las 580 horas propias del TFG y no las 980 horas
acumuladas durante todo el proyecto. La dedicación específica del TFG fue
superior a la carga académica teórica debido principalmente a la complejidad
de la integración con UVLHub y Flamapy, la adaptación de las pruebas, las
tareas adicionales de validación y las correcciones derivadas de las
revisiones técnicas.

\ph{Cierra la sección con una valoración global de la desviación temporal.}

\section{Plan de comunicaciones y seguimiento}
\label{sec:comunicaciones}

\ph{Describe cómo se coordina el trabajo con el tutor o tutores y cómo queda
constancia del avance. Tres filas bastan; lo que da valor es poder evidenciarlo
después (fechas de reuniones clave, decisiones tomadas en ellas, issues cerradas).}

\ph{Cierra con una o dos frases sobre cómo funcionó en la práctica: número
aproximado de reuniones mantenidas y un ejemplo concreto de decisión tomada en una
de ellas (con fecha).}

\begin{longtable}{@{}P{3.8cm} P{2.2cm} P{2.8cm} P{5.4cm}@{}}
\toprule
\textbf{Actividad} & \textbf{Cadencia} & \textbf{Vía} & \textbf{Resultado esperado} \\
\midrule
\endhead

Reuniones con el tutor &
Según avance del proyecto &
Presencial &
Revisión del estado del desarrollo, validación de avances, detección de posibles deficiencias y definición de las próximas tareas a realizar. \\[0.3em]

Mensajería asíncrona &
Bajo demanda &
Telegram y correo electrónico &
Resolución de dudas sobre el diseño y desarrollo del sistema, recepción de indicaciones adicionales, concertación de tutorías y evaluación de avances menores. \\[0.3em]

Seguimiento del desarrollo &
Continua &
Repositorio Git/GitHub &
Trazabilidad de la evolución del proyecto mediante el historial de cambios y commits realizados durante la implementación. \\

\bottomrule
\caption{Plan de comunicaciones y seguimiento del proyecto.}
\label{tab:comunicaciones}
\end{longtable}

Durante el desarrollo del proyecto se han mantenido diferentes reuniones
presenciales con el tutor con el objetivo de revisar el estado de la solución
y validar las decisiones tomadas durante la implementación. La primera reunión
relevante tuvo lugar el 5 de mayo de 2026, cuando se alcanzó aproximadamente
el 80\% del desarrollo. En ella se concluyó que la solución era
funcionalmente aceptable, aunque presentaba ciertas deficiencias técnicas que
debían ser corregidas.

Tras aplicar las correcciones necesarias, el 29 de junio de 2026 se realizó
una segunda revisión del estado del sistema y el 10 de julio se validó
positivamente la versión principal de la solución. Posteriormente, durante la
revisión técnica final, el tutor señaló nuevos aspectos relacionados con la
integración como \textit{plugin} de Flamapy, la validación de los modelos
generados y la organización de las pruebas. Estas observaciones dieron lugar
a las correcciones realizadas entre el 14 y el 18 de septiembre de 2026,
fecha en la que se considera cerrado técnicamente el desarrollo.

\section{Gestión de riesgos}
\label{sec:riesgos}

\ph{Identifica los riesgos principales que podían comprometer los objetivos o los
plazos. Para cada uno indica probabilidad, impacto y exposición (combinación de
ambos: Alta / Media / Baja), que ordena los riesgos por prioridad de vigilancia,
y el plan de mitigación. Regla para la exposición (matriz 3$\times$3 estándar de
probabilidad por impacto): exposición = el nivel más bajo de los dos, subiendo un
nivel si el impacto es Alto.}

\begin{longtable}{@{}P{0.7cm} P{3.2cm} P{1.4cm} P{1.8cm} P{2.4cm} P{4.4cm}@{}}
  \toprule
  \textbf{ID} & \textbf{Riesgo} & \textbf{Prob.} & \textbf{Impacto} & \textbf{Exposición} & \textbf{Mitigación} \\\\
  \midrule
  \endhead

  R1 &
  Complejidad de integración del generador dentro del ecosistema UVLHub-Flamapy &
  Alta &
  Alto &
  Alta &
  Realizar un estudio inicial de la arquitectura existente, validar progresivamente los cambios realizados y comprobar la compatibilidad entre los diferentes componentes del sistema. \\\\[0.3em]

  R2 &
  Problemas derivados de dependencias externas y configuración del entorno de ejecución &
  Media &
  Alto &
  Alta &
  Mantener actualizadas las dependencias necesarias, documentar la configuración del entorno y realizar pruebas de integración tras modificar versiones o paquetes externos. \\\\[0.3em]

  R3 &
  Errores en la generación de modelos tras la incorporación de nuevas funcionalidades &
  Alta &
  Medio &
  Media &
  Realizar validaciones funcionales tras cada modificación e incorporar pruebas automatizadas que permitan detectar errores en las funcionalidades implementadas. \\\\[0.3em]

  R4 &
  Subestimación del esfuerzo necesario para completar la implementación y validación del sistema &
  Media &
  Alto &
  Media &
  Realizar un seguimiento periódico del progreso, priorizar las tareas críticas y reservar tiempo adicional para corrección de errores y mejora de las pruebas. \\

  \bottomrule
  \caption{Riesgos identificados, exposición y planes de mitigación.}
  \label{tab:riesgos}
\end{longtable}

\subsection{Seguimiento de riesgos al cierre}
\label{sec:seguimiento_riesgos}

\ph{Un registro de riesgos solo demuestra gestión si se revisita al final. Indica,
para cada riesgo, si llegó a materializarse, qué respuesta se aplicó realmente
(la planificada u otra) y qué riesgo residual queda. Incluye una fila por cada
riesgo de la Tabla~\ref{tab:riesgos}, sin excepciones.}

\begin{longtable}{@{}P{0.7cm} P{2.8cm} P{5.4cm} P{5.3cm}@{}}
  \toprule
  \textbf{ID} & \textbf{¿Se materializó?} & \textbf{Respuesta aplicada} & \textbf{Riesgo residual} \\\\
  \midrule
  \endhead

  R1 &
  Sí &
  El riesgo apareció principalmente durante la fase inicial de comprensión de Flamapy y durante la integración del generador dentro de UVLHub. Para solucionarlo se estudió la arquitectura de Flamapy, se realizaron modificaciones sobre la estructura del código y se aplicó una estrategia de pruebas progresivas para validar cada incorporación realizada. &
  El generador se encuentra completamente integrado y cumple con la funcionalidad establecida inicialmente. No obstante, como ocurre en cualquier proyecto software, siempre existe margen de mejora relacionado con la reducción de deuda técnica y la optimización de la arquitectura interna. \\\\[0.3em]

  R2 &
  Sí &
  Se produjeron problemas relacionados con las dependencias externas del proyecto, especialmente con la actualización de Flamapy, la creación y configuración de los wheels necesarios para Pyodide y los problemas encontrados para que estos paquetes fueran localizados y utilizados correctamente. La solución consistió en actualizar dependencias, crear nuevos wheels cuando fue necesario e investigar la configuración requerida para su correcta integración. &
  El sistema funciona correctamente con las versiones actuales de las dependencias utilizadas. Sin embargo, futuras actualizaciones de estas tecnologías podrían requerir nuevos ajustes para mantener la compatibilidad. \\\\[0.3em]

  R3 &
  Sí &
  Durante el desarrollo se detectaron en varias ocasiones problemas en la generación de modelos debido a que determinadas modificaciones no producían los resultados esperados. Para solucionarlos se comprobaba inicialmente el fallo mediante la generación de diferentes modelos, posteriormente se analizaba la causa, se corregía la implementación y se realizaban nuevas validaciones manuales, actualizando las pruebas unitarias cuando aportaban valor. &
  La generación de modelos ha sido validada de forma exhaustiva siguiendo los criterios establecidos. Aun así, no se descarta la aparición de nuevos casos límite ante futuras ampliaciones o nuevos escenarios de uso. \\\\[0.3em]

  R4 &
  Sí &
  La estimación inicial del esfuerzo necesario se vio superada debido a diversos factores relacionados con la calidad final del producto, como los errores detectados durante la implementación, la necesidad de mejorar el enfoque de las pruebas, los problemas de integración y las tareas de refactorización necesarias. La gestión de este riesgo se realizó priorizando las correcciones necesarias, reorganizando tareas y utilizando las revisiones con el tutor para obtener feedback sobre el estado del proyecto. &
  El principal riesgo residual está relacionado con la posible existencia de deuda técnica pendiente. Una mayor reducción de esta deuda mediante refactorizaciones y mejoras arquitectónicas podría facilitar futuras ampliaciones y reducir el esfuerzo necesario para incorporarlas. \\

  \bottomrule
  \caption{Seguimiento de riesgos al cierre: materialización, respuesta y residuales.}
  \label{tab:seguimiento_riesgos}
\end{longtable}

\section{Control de cambios}
\label{sec:control_cambios}

\ph{El plan es un documento vivo: registra los cambios sobre la línea base
(alcance, plazo o enfoque) acordados durante el proyecto. Anti-burocracia: recoge
solo los 3--6 cambios SIGNIFICATIVOS respecto a la línea base; no registres
microajustes. Ten en cuenta que un TFG sin ningún cambio registrado resulta poco
creíble. Si un cambio afectó al alcance, verifica que la Sección~\ref{sec:alcance}
quedó coherente.}


\begin{longtable}{@{}P{0.6cm} P{2.2cm} P{3.1cm} P{2.5cm} P{3.2cm} P{2.5cm}@{}}
  \toprule
  \textbf{ID} & \textbf{Fecha} & \textbf{Cambio} & \textbf{Causa} & 
  \textbf{Impacto en la línea base} & \textbf{Aprobado por} \\\\
  \midrule
  \endhead

  C1 &
  12/05/2025 &
  Ampliación de las opciones de configuración del generador para aumentar la flexibilidad del proceso de generación. &
  Necesidad de proporcionar una mayor capacidad de adaptación de los modelos generados a distintos escenarios de uso. &
  Incremento del alcance funcional mediante la incorporación de nuevas opciones de configuración relacionadas con la generación de operadores y atributos. &
  Tutor / investigadores \\\\[0.8em]

  C2 &
  Durante el desarrollo &
  Ampliación del plazo previsto para completar la implementación y validación del generador automático. &
  La estimación inicial no contemplaba completamente la complejidad de la integración con el ecosistema UVLHub-Flamapy, así como los problemas surgidos durante la implementación y corrección de errores. &
  Desplazamiento del hito previsto de finalización del desarrollo y aumento del esfuerzo dedicado a la implementación y validación del sistema. &
  Tutor \\\\[0.8em]

  C3 &
  05/05/2026 &
  Revisión del enfoque técnico y refactorización de la estructura interna del generador para mejorar su integración y mantenibilidad. &
  Durante la revisión del estado del proyecto se identificaron deficiencias técnicas relacionadas con la organización interna del código y la integración del generador dentro del ecosistema existente. &
  Necesidad de reorganizar componentes internos, mejorar la arquitectura de la solución y adaptar parte de la implementación existente. &
  Tutor, reunión de seguimiento 05/05/2026 \\\\[0.8em]

  C4 &
  05/2026 &
  Sustitución del front-end inicial por la propuesta proporcionada por el tutor para la integración del generador en UVLHub. &
  La propuesta alternativa ofrecía una mejor adaptación al flujo existente de la plataforma y una experiencia de usuario más adecuada para configurar el proceso de generación. &
  Modificación de la implementación del front-end y adaptación del generador al nuevo flujo de interacción, manteniendo los objetivos funcionales definidos inicialmente. &
  Tutor \\\\[0.8em]

  C5 &
  14--18/09/2026 &
  Corrección y consolidación final del generador tras la revisión técnica del
  tutor. &
  Observaciones relativas a la integración como \textit{plugin} de Flamapy,
  la validación de modelos UVL, la generación satisfacible y la organización
  de las pruebas. &
  Extensión del cierre técnico hasta septiembre de 2026, sin modificar el
  alcance funcional del proyecto. &
  Tutor \\  

  \bottomrule
  \caption{Registro de control de cambios sobre la línea base.}
  \label{tab:control_cambios}
\end{longtable}

\section{Plan de calidad}
\label{sec:plan_calidad}

\ph{Declara brevemente qué criterios de calidad se exigen al producto y a la
memoria, y con qué mecanismo se verifica cada uno. No dupliques otros capítulos:
referencia. Media página como máximo. No lo confundas con la metodología de
validación (Sección~\ref{sec:metodologia_validacion}): aquella define CÓMO se mide
el cumplimiento de los objetivos; el plan de calidad fija los estándares del
proceso y del producto (código, revisiones, umbrales) con los que se trabaja a
diario.}

\ph{\begin{itemize}
  \item \ph{\textbf{Código:} estándares de estilo y \textit{linting} automatizado; cómo se aplican se detalla en el Capítulo~\ref{chap:implementacion}.}
  \item \ph{\textbf{Entregables:} criterio de aceptación de cada paquete de la EDT, verificado mediante las pruebas del Capítulo~\ref{chap:pruebas}.}
  \item \ph{\textbf{Cobertura:} umbral objetivo de cobertura de pruebas (p.\ ej.\ X\,\%), medido como se describe en la Sección~\ref{sec:cobertura}.}
  \item \ph{\textbf{Memoria:} revisión de cada capítulo por el tutor/a y verificación de referencias cruzadas y bibliografía antes de la entrega.}
\end{itemize}}

El plan de calidad definido para el proyecto establece los criterios utilizados para garantizar la calidad tanto de la solución software desarrollada como de la documentación asociada al TFG.

\begin{itemize}

\item \textbf{Código:} se ha buscado mantener una implementación organizada, legible y mantenible, siguiendo una estructura coherente con la arquitectura definida para el generador. Para facilitar la comprensión del código y detectar posibles problemas de estilo, se ha utilizado una herramienta de \textit{linting} integrada dentro del flujo de desarrollo mediante GitHub Actions, ya incorporada en la configuración de UVLHub. Además, durante el desarrollo se realizaron tareas de refactorización orientadas a mejorar la calidad interna del código y reducir posibles problemas derivados de la evolución del sistema. La trazabilidad de los cambios se ha mantenido mediante el control de versiones con Git y el uso de commits siguiendo la convención \textit{Conventional Commits}. La siguiente imagen muestra los workflows usados:

\begin{figure}[H]
    \centering
    \includegraphics[width=0.5\textwidth]{workflows.png}
    \caption{Arquitectura general de UVLHub y sus integraciones principales. Elaboración propia.}
    \label{fig:workflows}
\end{figure}

\item \textbf{Entregables:} cada bloque funcional desarrollado debía cumplir con los objetivos establecidos y aportar una funcionalidad verificable dentro del sistema. Para ello, las nuevas incorporaciones se validaban progresivamente antes de continuar con el desarrollo de nuevas capacidades, asegurando que los componentes implementados mantenían el comportamiento esperado.

\item \textbf{Pruebas:} la calidad funcional del sistema se ha comprobado mediante la combinación de validaciones manuales y pruebas automatizadas. Tras la incorporación de nuevas funcionalidades se verificaba que los modelos generados cumplían los criterios esperados y que la interacción mediante el front-end funcionaba correctamente. Además, se incorporaron pruebas unitarias adicionales cuando estas aportaban valor al detectar posibles errores relevantes en el sistema.

\item \textbf{Memoria:} la calidad de la documentación se ha garantizado mediante revisiones periódicas de la estructura y contenido de la memoria, así como mediante la revisión de los avances con el tutor del proyecto. Antes de la entrega final se realizará una comprobación de la coherencia entre los diferentes capítulos, referencias cruzadas y bibliografía utilizada.

\end{itemize}

\section{Presupuesto}
\label{sec:presupuesto}

\ph{Presenta el presupuesto del proyecto. Una opción habitual es seguir el modelo
TCO (\textit{Total Cost of Ownership}), separando costes de capital (CapEx) y costes
operacionales (OpEx). Conviene incluir, junto al coste real (a menudo 0\,€ en contexto
académico), el valor de mercado equivalente para una valoración realista.}

\subsection*{CapEx --- Costes de capital}

\ph{Inversiones en infraestructura física y licencias de uso prolongado.}

\ph{Justifica brevemente cómo se ha estimado el valor de mercado de cada recurso
(por ejemplo, tarifa pública de un proveedor cloud o precio comercial de una licencia).}

\begin{longtable}{P{4cm} P{4cm} P{2.5cm} P{2.5cm}}
  \toprule
  \textbf{Recurso} & \textbf{Descripción} & \textbf{Coste real} & \textbf{Valor} \\
  \midrule

  Equipo de desarrollo &
  Equipo personal utilizado para la implementación, pruebas y mantenimiento del sistema. &
  0,00\,€ &
  X\,€ \\[0.3em]

  Licencias software &
  Herramientas de desarrollo y tecnologías empleadas durante el proyecto, basadas principalmente en software libre. &
  0,00\,€ &
  X\,€ \\

  \midrule

  \textbf{Total CapEx} &
  &
  \textbf{0,00\,€} &
  \textbf{X\,€} \\

  \bottomrule
  \caption{Costes de capital (CapEx).}
  \label{tab:capex}
\end{longtable}

Los recursos utilizados durante el desarrollo no supusieron una inversión económica directa conocida, ya que se empleó principalmente infraestructura propia y herramientas basadas en software libre o de coste nulo. Por tanto, el coste real asociado a los costes de capital del proyecto es de 0,00\,€.

\subsection*{OpEx --- Costes operacionales}

\ph{Gastos recurrentes asociados al desarrollo y mantenimiento.}

\subsubsection*{Personal}

\ph{Indica el rol, la tarifa por hora y la justificación de esa tarifa
(convenio colectivo, mercado laboral de referencia, etc.). Usa como actividades
las fases de la Tabla~\ref{tab:resumen-global-proyecto}; los totales de horas deben
coincidir con ella.}


El coste de personal corresponde a la dedicación asociada al contrato de
apoyo técnico vinculado al proyecto TASOVA Plus. Dicho contrato tuvo una
duración de cuatro meses, entre abril y julio de 2025, y contemplaba una
dedicación de 25 horas semanales junto con una remuneración bruta mensual
de 930,60\,€.

La tarifa utilizada para calcular el coste estimado se ha fijado de forma
uniforme en 9,31\,€ por hora y se aplica a las 750 horas de desarrollo
previstas inicialmente. Por este motivo, los costes estimados de la tabla
se obtienen multiplicando las horas estimadas de cada fase por dicha tarifa.

El coste real refleja únicamente la remuneración correspondiente al
contrato de TASOVA Plus. Las cuatro mensualidades abonadas durante el
periodo contractual suponen un coste real total de
\(4 \times 930,60 = 3722,40\,€\), equivalente aproximadamente a
9,306\,€ por hora para las 400 horas contractuales. Las horas realizadas
específicamente para el TFG se desarrollaron fuera del periodo financiado y
no representan un coste directo adicional de personal.

\begin{table}[H]
\centering
\small
\setlength{\tabcolsep}{3pt}
\renewcommand{\arraystretch}{1.25}

\begin{tabularx}{\textwidth}{@{}
    >{\raggedright\arraybackslash}X
    >{\centering\arraybackslash}p{1.35cm}
    >{\centering\arraybackslash}p{1.45cm}
    >{\centering\arraybackslash}p{1.15cm}
    >{\centering\arraybackslash}p{1.25cm}
    >{\centering\arraybackslash}p{2.05cm}
    >{\centering\arraybackslash}p{1.80cm}
@{}}
\toprule

\textbf{Actividad} &
\shortstack[c]{\textbf{H.}\\\textbf{esti-}\\\textbf{madas}} &
\shortstack[c]{\textbf{H.}\\\textbf{TASO-}\\\textbf{VA}} &
\shortstack[c]{\textbf{H.}\\\textbf{TFG}} &
\shortstack[c]{\textbf{H.}\\\textbf{reales}} &
\shortstack[c]{\textbf{Coste}\\\textbf{estimado}\\\textbf{(€)}} &
\shortstack[c]{\textbf{Coste}\\\textbf{real}\\\textbf{(€)}} \\

\midrule

Análisis inicial y planificación del proyecto &
12 & 12 & 0 & 12 & 111,72 & 111,67 \\

Análisis y diseño de la solución &
25 & 25 & 0 & 25 & 232,75 & 232,65 \\

Desarrollo e integración del generador automático &
593 & 363 & 340 & 703 & 5520,83 & 3378,08 \\

Pruebas y validación de la solución &
120 & 0 & 170 & 170 & 1117,20 & 0,00 \\

\midrule

\textbf{Total personal} &
\textbf{750} &
\textbf{400} &
\textbf{510} &
\textbf{910} &
\textbf{6982,50} &
\textbf{3722,40} \\

\bottomrule
\end{tabularx}

\caption{Desglose de horas y costes de personal.}
\label{tab:coste_personal}
\end{table}

\subsubsection*{Herramientas y servicios de desarrollo}

\begin{longtable}{P{4cm} P{4cm} P{1.8cm} P{2.7cm}}
  \toprule
  \textbf{Herramienta} &
  \textbf{Licencia utilizada} &
  \textbf{Coste real} &
  \textbf{Valor de mercado} \\

  \midrule

  GitHub &
  Plataforma utilizada para almacenar el código fuente y gestionar el control de versiones. &
  0\,€ &
  0\,€ \\[0.3em]

  Render &
  Servicio utilizado para desplegar una versión accesible de UVLHub durante las revisiones funcionales del sistema. &
  0\,€ &
  X\,€ / mes \\[0.3em]

  Python, Flask, Flamapy y librerías asociadas &
  Tecnologías y dependencias utilizadas durante el desarrollo bajo licencias abiertas. &
  0\,€ &
  0\,€ \\

  \midrule

  \textbf{Total herramientas} &
  &
  \textbf{0\,€} &
  \textbf{X\,€} \\

  \bottomrule
  \caption{Costes de herramientas y servicios de desarrollo (OpEx).}
  \label{tab:coste_herramientas}
\end{longtable}

\subsection*{Resumen del presupuesto (TCO)}

\ph{Cierra la sección comentando el coste real frente al valor de mercado
y la influencia del entorno académico (licencias educativas, equipos cedidos\dots).}

\begin{longtable}{P{4cm} P{3cm} P{3cm} P{3cm}}
  \toprule
  \textbf{Categoría} &
  \textbf{Tipo} &
  \textbf{Coste real} &
  \textbf{Valor de mercado} \\

  \midrule

  Infraestructura y licencias &
  CapEx &
  0,00\,€ &
  X\,€ \\

  Personal &
  OpEx &
  3722,40\,€ &
  X\,€ \\

  Herramientas y servicios &
  OpEx &
  0,00\,€ &
  X\,€ \\

  \midrule

  \textbf{TCO total} &
  &
  \textbf{3722,40\,€} &
  \textbf{X\,€} \\

  \bottomrule
  \caption{Resumen del presupuesto del proyecto (TCO = CapEx + OpEx).}
  \label{tab:presupuesto_total}
\end{longtable}

El coste real del proyecto está determinado principalmente por los costes de personal asociados al contrato de apoyo técnico vinculado al proyecto TASOVA PLUS. El resto de recursos empleados no han supuesto costes adicionales debido al uso de herramientas de código abierto y servicios con planes gratuitos. Cabe destacar que, aunque el desarrollo continuó posteriormente fuera del periodo financiado, dicho esfuerzo adicional no representa un coste económico directo dentro del presupuesto del proyecto.

\section{Cierre del proyecto y lecciones aprendidas}
\label{sec:cierre}

\ph{Cierra formalmente el plan. Primero, verifica el proyecto contra la EDT
(Tabla~\ref{tab:edt}): todo paquete de trabajo debe estar entregado o tener un
cambio registrado en la Sección~\ref{sec:control_cambios} que lo descarte. Después,
resume en una frase la desviación global (Tabla~\ref{tab:resumen-global-proyecto}) y extrae
3--5 lecciones aprendidas honestas: las más valiosas son las que admiten un error
y explican qué se haría distinto.}

\ph{Las conclusiones (Capítulo~\ref{chap:conclusiones}) retomarán esta reflexión
de forma resumida: no dupliques aquí lo que dirás allí, ni allí lo que ya queda
registrado aquí.}


Al finalizar el proyecto, todos los paquetes de trabajo definidos en la EDT han sido completados. Las fases planificadas se han ejecutado en su totalidad, sin necesidad de descartar ningún paquete de trabajo. Los cambios relevantes realizados durante el desarrollo han quedado registrados en la Sección~\ref{sec:control_cambios}, manteniendo la coherencia con el alcance inicial definido.

La desviación temporal global del proyecto se ha concentrado principalmente en las fases de desarrollo e integración del generador automático y pruebas y validación de la solución, debido a la complejidad añadida durante la integración con componentes externos y a la necesidad de realizar iteraciones adicionales para asegurar el correcto funcionamiento del sistema.

\begin{longtable}{P{6cm} P{6cm}}
  \toprule
  \textbf{Lección} & \textbf{Qué se haría distinto} \\
  
  \midrule
  
  La estimación inicial del esfuerzo necesario para las fases de integración y validación fue inferior al esfuerzo realmente requerido. &
  Realizar una estimación más conservadora, incluyendo un margen adicional para problemas derivados de dependencias externas y correcciones durante la fase final. \\[0.5em]
  
  El conocimiento inicial sobre algunas tecnologías utilizadas en el proyecto no era suficiente para prever completamente las dificultades de integración. &
  Dedicar una fase inicial más extensa al estudio de las herramientas y componentes externos antes de comenzar la implementación. \\[0.5em]
  
  La estrategia de pruebas podría haberse definido con mayor profundidad desde las primeras etapas del desarrollo. &
  Diseñar desde el inicio una planificación de pruebas más completa, identificando las funcionalidades críticas y los casos de validación necesarios. \\[0.5em]
  
  La acumulación progresiva de cambios durante el desarrollo generó cierta deuda técnica. &
  Realizar revisiones y refactorizaciones periódicas del código para mantener una arquitectura más mantenible durante todo el ciclo de desarrollo. \\
  
  \bottomrule
  \caption{Lecciones aprendidas al cierre del proyecto.}
  \label{tab:lecciones-aprendidas}
\end{longtable}

En conjunto, el proyecto se ha completado con una dedicación real de
980 horas, de las cuales 910 corresponden al desarrollo y 70 a la elaboración
de la memoria. Aunque la planificación inicial se vio afectada por un periodo
de inactividad y por varias iteraciones adicionales de integración y
validación, todos los paquetes de trabajo definidos en la EDT fueron
completados.

Con el proyecto planificado, ejecutado y cerrado, el siguiente capítulo
precisa qué debe hacer exactamente la solución desarrollada. El
Capítulo~\ref{chap:requisitos} presenta el catálogo de requisitos funcionales
y no funcionales que sirve como base para el diseño, la implementación y la
validación del sistema.

% --- Transición al siguiente capítulo ---
\ph{Cierra el capítulo con un párrafo puente de 2--4 líneas: qué queda
establecido aquí y qué pregunta abre el capítulo siguiente. Por ejemplo:
«Con el plan establecido (qué se hará, cuándo y con qué control), el siguiente
paso es precisar QUÉ debe hacer exactamente el sistema: el
Capítulo~\ref{chap:requisitos} traduce los objetivos en un catálogo de requisitos
verificables».}